\documentclass[12pt,a4paper]{article}

\usepackage[margin=2.5cm]{geometry}
\usepackage{graphicx}
\usepackage{booktabs}
\usepackage{amsmath}
\usepackage{caption}
\usepackage{subcaption}
\usepackage{float}
\usepackage{hyperref}
\usepackage{parskip}
\usepackage{titlesec}
\usepackage{xcolor}

\hypersetup{colorlinks=true, linkcolor=blue, urlcolor=blue}

\title{\textbf{Assignment \#2: Heart Failure Classification}\\
\large CSE: Pattern Recognition\\
Alexandria University -- Faculty of Engineering}
\author{
  Nour Eldin Akram  	\quad \texttt{22011309} \\
  Ahmed Ragy Hussein	 \quad \texttt{22011690} \\
  Mina Nabil Youssef	 \quad \texttt{22011290} \\[6pt]
  \normalsize Computer and Systems Engineering Department
}
\date{April 2026}

\begin{document}

\maketitle
\tableofcontents
\newpage

%=============================================================
\section{Problem Statement}
%=============================================================

Heart failure is a critical condition caused by cardiovascular diseases. The goal of this
assignment is to build and compare several classifiers to predict whether a patient is at
risk of heart failure based on clinical features. The dataset used is the
\textit{Heart Failure Prediction} dataset (918 samples, 11 features, 2 classes) from Kaggle.

%=============================================================
\section{Dataset and Preparation}
%=============================================================

The dataset was prepared following the steps below:

\begin{itemize}
  \item \textbf{Random seed:} Fixed at 42 for all random operations to ensure reproducibility.
  \item \textbf{Train/Validation/Test split:} 70\% training, 10\% validation, 20\% test,
        using stratified splitting to preserve class distribution.
  \item \textbf{Feature preprocessing:} Numerical features were standardized.
        Categorical features (e.g.\ \texttt{Sex}, \texttt{ChestPainType}) were one-hot encoded.
  \item \textbf{Validation set role:} Used exclusively for hyperparameter tuning and model
        selection; the test set was kept unseen until final evaluation.
\end{itemize}

%=============================================================
\section{Implementation Details}
%=============================================================

Three classifiers were implemented from scratch, plus a bonus Random Forest model.

\subsection{Decision Tree}
A decision tree was implemented using Information Gain (entropy-based) as the splitting
criterion. Two hyperparameters were tuned on the validation set via a grid search:
\texttt{max\_depth} $\in \{2,3,\ldots,30\}$ and
\texttt{min\_samples\_split} $\in \{2,5,7,10,20,30,40\}$.

\subsection{Bagging Ensemble}
A Bagging (Bootstrap Aggregating) classifier was built using the custom decision tree as
the base learner. Multiple trees are trained on bootstrap samples of the training data;
predictions are aggregated by majority vote.
The number of estimators \texttt{n\_estimators} $\in \{1,5,10,\ldots,100\}$ was tuned on
the validation set.

\subsection{AdaBoost Ensemble}
AdaBoost was implemented from scratch with depth-1 decision stumps as weak learners.
Sample weights are updated after each round so that subsequent stumps focus on previously
misclassified examples. The number of boosting iterations was tuned on the validation set
across $\{1,5,10,20,30,50,75,100,150,200\}$.

\subsection{Random Forest (Bonus)}
A Random Forest was implemented, extending the Bagging approach by additionally applying
random feature subsampling at each split. The number of estimators was tuned identically
to the Bagging setup.

%=============================================================
\section{Hyperparameter Tuning}
%=============================================================

\subsection{Decision Tree}

Figure~\ref{fig:dt_tuning} shows validation accuracy and F1-score across the
(\texttt{max\_depth}, \texttt{min\_samples\_split}) grid. Shallow trees (depth 2--5)
with larger \texttt{min\_samples\_split} values achieve the highest validation scores.
The best configuration found was \textbf{max\_depth = 5, min\_samples\_split = 2},
yielding a validation accuracy of 0.8791 and F1-score of 0.8889.

\begin{figure}[H]
  \centering
  \begin{subfigure}[b]{0.48\textwidth}
    \includegraphics[width=\textwidth]{images/dt_acc_tuning.png}
    \caption{Validation Accuracy}
  \end{subfigure}
  \hfill
  \begin{subfigure}[b]{0.48\textwidth}
    \includegraphics[width=\textwidth]{images/dt_f1_tuning.png}
    \caption{Validation F1-Score}
  \end{subfigure}
  \caption{Decision Tree hyperparameter tuning heatmaps.}
  \label{fig:dt_tuning}
\end{figure}

\subsection{AdaBoost}

Figure~\ref{fig:ada_tuning} shows that performance rises sharply in the first 50 iterations,
peaks at \textbf{75 iterations} (validation accuracy 0.9121, F1 0.9216), and then slightly
decreases due to overfitting at higher iteration counts.

\begin{figure}[H]
  \centering
  \begin{subfigure}[b]{0.48\textwidth}
    \includegraphics[width=\textwidth]{images/ada_acc.png}
    \caption{Validation Accuracy vs.\ Iterations}
  \end{subfigure}
  \hfill
  \begin{subfigure}[b]{0.48\textwidth}
    \includegraphics[width=\textwidth]{images/ada_f1.png}
    \caption{Validation F1-Score vs.\ Iterations}
  \end{subfigure}
  \caption{AdaBoost tuning curves.}
  \label{fig:ada_tuning}
\end{figure}

\subsection{Bagging}

Figure~\ref{fig:bag_tuning} shows that accuracy and F1-score plateau around
\texttt{n\_estimators = 20--30}. The best model used \textbf{n\_estimators = 20}
(validation accuracy 0.8901, F1 0.9000).

\begin{figure}[H]
  \centering
  \begin{subfigure}[b]{0.48\textwidth}
    \includegraphics[width=\textwidth]{images/bag_acc.png}
    \caption{Validation Accuracy vs.\ Trees}
  \end{subfigure}
  \hfill
  \begin{subfigure}[b]{0.48\textwidth}
    \includegraphics[width=\textwidth]{images/bag_f1.png}
    \caption{Validation F1-Score vs.\ Trees}
  \end{subfigure}
  \caption{Bagging tuning curves.}
  \label{fig:bag_tuning}
\end{figure}

\subsection{Random Forest (Bonus)}

Figure~\ref{fig:rf_tuning} shows a sharp peak at \textbf{n\_estimators = 20}
(validation accuracy 0.9341, F1 0.9423), which is the highest validation performance
across all models. Performance degrades slightly as more trees are added.

\begin{figure}[H]
  \centering
  \begin{subfigure}[b]{0.48\textwidth}
    \includegraphics[width=\textwidth]{images/rf_acc.png}
    \caption{Validation Accuracy vs.\ Trees}
  \end{subfigure}
  \hfill
  \begin{subfigure}[b]{0.48\textwidth}
    \includegraphics[width=\textwidth]{images/rf_f1.png}
    \caption{Validation F1-Score vs.\ Trees}
  \end{subfigure}
  \caption{Random Forest tuning curves.}
  \label{fig:rf_tuning}
\end{figure}

%=============================================================
\section{Evaluation}
%=============================================================

\subsection{Performance Summary}

Table~\ref{tab:results} summarises the test-set performance of all models after
hyperparameter selection on the validation set.

\begin{table}[H]
  \centering
  \caption{Test-set performance comparison of all classifiers.}
  \label{tab:results}
  \begin{tabular}{lcc}
    \toprule
    \textbf{Model} & \textbf{Test Accuracy} & \textbf{Test F1-Score} \\
    \midrule
    Decision Tree   & 0.8324 & 0.8442 \\
    Bagging         & 0.8432 & 0.8585 \\
    AdaBoost        & 0.8649 & 0.8744 \\
    Random Forest (Bonus) & \textbf{0.8865} & \textbf{0.8966} \\
    \bottomrule
  \end{tabular}
\end{table}

\subsection{Confusion Matrices}

\subsubsection{Decision Tree}

\begin{figure}[H]
  \centering
  \begin{subfigure}[b]{0.45\textwidth}
    \includegraphics[width=\textwidth]{images/dt_val_cm.png}
    \caption{Validation set}
  \end{subfigure}
  \hfill
  \begin{subfigure}[b]{0.45\textwidth}
    \includegraphics[width=\textwidth]{images/dt_test_cm.png}
    \caption{Test set}
  \end{subfigure}
  \caption{Decision Tree confusion matrices (classes: 0 = No Heart Failure, 1 = Heart Failure).}
  \label{fig:dt_cm}
\end{figure}

On the test set the Decision Tree correctly classifies 70 true-negatives and
84 true-positives, while misclassifying 12 negatives as positive and
19 positives as negative.

\subsubsection{Bagging}

\begin{figure}[H]
  \centering
  \begin{subfigure}[b]{0.45\textwidth}
    \includegraphics[width=\textwidth]{images/bag_val_cm.png}
    \caption{Validation set}
  \end{subfigure}
  \hfill
  \begin{subfigure}[b]{0.45\textwidth}
    \includegraphics[width=\textwidth]{images/bag_test_cm.png}
    \caption{Test set}
  \end{subfigure}
  \caption{Bagging confusion matrices.}
  \label{fig:bag_cm}
\end{figure}

The Bagging ensemble achieves 68 true-negatives and 88 true-positives on the test set,
with 14 false-positives and 15 false-negatives.

\subsubsection{AdaBoost}

\begin{figure}[H]
  \centering
  \begin{subfigure}[b]{0.45\textwidth}
    \includegraphics[width=\textwidth]{images/ada_val_cm.png}
    \caption{Validation set}
  \end{subfigure}
  \hfill
  \begin{subfigure}[b]{0.45\textwidth}
    \includegraphics[width=\textwidth]{images/ada_test_cm.png}
    \caption{Test set}
  \end{subfigure}
  \caption{AdaBoost confusion matrices (classes: $-1$ = No Heart Failure, $+1$ = Heart Failure).}
  \label{fig:ada_cm}
\end{figure}

AdaBoost achieves 73 true-negatives and 87 true-positives on the test set,
with only 9 false-positives and 16 false-negatives.

\subsubsection{Random Forest (Bonus)}

\begin{figure}[H]
  \centering
  \begin{subfigure}[b]{0.45\textwidth}
    \includegraphics[width=\textwidth]{images/rf_val_cm.png}
    \caption{Validation set}
  \end{subfigure}
  \hfill
  \begin{subfigure}[b]{0.45\textwidth}
    \includegraphics[width=\textwidth]{images/rf_test_cm.png}
    \caption{Test set}
  \end{subfigure}
  \caption{Random Forest confusion matrices.}
  \label{fig:rf_cm}
\end{figure}

The Random Forest achieves the best test results: 73 true-negatives and
91 true-positives, with only 9 false-positives and 12 false-negatives.

%=============================================================
\section{Analysis and Comparison}
%=============================================================

\subsection{Which Model Performed Best and Why?}

The \textbf{Random Forest} achieved the highest test accuracy (0.8865) and F1-score
(0.8966), followed by AdaBoost (accuracy 0.8649, F1 0.8744). Both ensemble methods
significantly outperform the single Decision Tree (accuracy 0.8324).

Random Forest benefits from two sources of variance reduction: bootstrap sampling
(as in Bagging) \emph{and} random feature subsampling at each split. This decorrelates
the individual trees more aggressively than Bagging alone, leading to a more powerful
ensemble.

\subsection{Decision Tree vs.\ Ensemble Methods}

A single Decision Tree is prone to overfitting: it perfectly memorises the training data
unless pruned, yet generalises less well. The tuned tree (depth 5) achieved only 0.8324
on the test set. Ensemble methods address this by averaging over many trees trained on
different data subsets (Bagging, Random Forest) or by iteratively correcting errors
(AdaBoost), all of which reduce variance or bias.

\begin{itemize}
  \item \textbf{Bagging} reduces variance by averaging predictions, achieving 0.8432 test accuracy.
  \item \textbf{AdaBoost} reduces bias by focusing on hard examples; it scores 0.8649,
        clearly outperforming the single tree.
  \item \textbf{Random Forest} combines variance reduction with decorrelated trees,
        achieving the best overall performance.
\end{itemize}

\subsection{Misclassification Patterns}

Across all models, false-negatives (heart failure patients predicted as healthy) are more
common than false-positives. In the medical context of heart failure prediction, false-negatives
are particularly costly because a missed diagnosis can have severe clinical consequences.
AdaBoost shows the fewest false-positives (9) on the test set, while the Decision Tree has
the most false-negatives (19).

The Random Forest achieves the best balance: only 9 false-positives and 12 false-negatives,
making it the most reliable model for this medical classification task.

%=============================================================
\section{Conclusion}
%=============================================================

This assignment compared four classifiers -- Decision Tree, Bagging, AdaBoost, and
Random Forest -- on the Heart Failure Prediction dataset. All three ensemble methods
outperformed the single decision tree, confirming the well-known advantages of ensemble
learning. The Random Forest achieved the best test-set performance (accuracy 0.8865,
F1-score 0.8966), making it the recommended model for this task.

\end{document}